\documentclass[11pt]{article}
\usepackage[margin=2.5cm]{geometry}
\usepackage{iftex}
\ifPDFTeX
  \usepackage[T1]{fontenc}
  \usepackage{lmodern}
  \usepackage{textcomp}
  \DeclareUnicodeCharacter{2010}{-}
  \DeclareUnicodeCharacter{2011}{-}
\else
  \usepackage{fontspec}
\fi
\usepackage{graphicx}
\usepackage[numbers,sort&compress]{natbib}
\usepackage{amsmath,amssymb}
\usepackage{booktabs,longtable,array,calc}
\usepackage[hidelinks]{hyperref}
\usepackage{xcolor}
\usepackage{lineno}
\linenumbers
\providecommand{\tightlist}{\setlength{\itemsep}{0pt}\setlength{\parskip}{0pt}}
\newcommand{\passthrough}[1]{#1}
\setlength{\parskip}{4pt}
\setlength{\emergencystretch}{3em}\sloppy
\title{Supplementary Materials for: The fault architecture and seismic behavior of an entire plate-boundary system offshore Cascadia}
\author{[Authors and affiliations to add]}
\date{Draft, generated from \texttt{paper\_sciadv/supplementary.qmd} in Denolle-Lab/cascadia\_obs\_ensemble}
\begin{document}
\maketitle
\section{Text S1. Detecting fault planes in
seismicity}\label{text-s1.-detecting-fault-planes-in-seismicity}

\textbf{Why planes, not clusters.} A first version found clusters by
density (DBSCAN: eight neighbors within 2.5 km in 3D). It kept only the
densest knots (20\% of events) and missed long, sparse alignments such
as the plate boundaries, because density never tests whether events lie
on a plane. The method used here searches for planes directly, in the
spirit of anisotropic clustering \citep{Ouillon2008, Wang2013} and of
recent hypocenter-based fault models \citep{Alongi2025}.

\textbf{Detection.} (Code:
\texttt{workflow/07\_faults/fault\_surfaces\_planes.py}.) - Events: at
least four P and four S picks, catalog depth \ensuremath{\geq} 0, Axial Seamount
excluded (76,281 events). - Tolerance:
\(\tau = 1.5\sqrt{(1.5\,|n_h|)^2 + (3\,n_z)^2}\) km for a plane with
unit normal \(n\) (about 2.3 km for a vertical plane, 4.5 km for a
horizontal one). - Seeds from dense to sparse; 150 random planes per
seed through the seed and two events within 10 km, scored in a 20 km
ball \citep{Fischler1981}; growth through events within \(\tau\) that
are connected in the plane (4 km cells, 8-connectivity); acceptance with
\ensuremath{\geq} 30 events, planarity \(\sqrt{\lambda_3/\lambda_2} \le
  0.5\), in-plane width \ensuremath{\geq} 1 km, dip \ensuremath{\geq} 30\textdegree{} (v4 depths cannot resolve
flatter structures; a diffuse cloud passes for a flat plane). -
Reassignment (three passes, k-planes): each event to the plane with the
smallest \(|d|/\tau\) whose rectangle contains its projection. -
Merging: normals within 15\textdegree{}, each within \(\tau\) of the other, events
closer than 4 km, joined set still planar. - Rectangles: the 2--98\%
range of the events along strike and down dip plus 0.5 km; plane or
quadratic by BIC; 1 km grid.

\textbf{Result.} 155 faults detected; 10 removed (Text S3); Blanco
redone as strands (Text S2); 141 in the product, 70\% of events assigned
(Table S1: \texttt{data/faults/v0/fault\_surfaces\_v0\_assignment.csv}).
Fig. S1: map of events colored by fault
(\texttt{fault\_surfaces\_v0\_map.png}).

\section{Text S2. Uncertainty, stability and Blanco
strands}\label{text-s2.-uncertainty-stability-and-blanco-strands}

\textbf{Uncertainty per vertex} (1\ensuremath{\sigma} normal to the surface): bootstrap of
the fault's events (100 draws); Monte Carlo moving each event by N(0;
1.5, 1.5, 3 km); spread over model choices (plane or quadratic, with or
without trimming beyond 2.5 RMS); combined in quadrature, median 0.44 km
(10th--90th percentile 0.17--0.78 km). The catalog's
\texttt{spatial\_sigma} (3--20 km) is not used: events scatter only
0.6--1.8 km about the fitted surfaces, so it behaves like an absolute,
not a relative, error. Strike and dip intervals come from the bootstrap;
the full ensemble is provided
(\texttt{fault\_surfaces\_v0\_ensemble.csv.gz}).

\textbf{Stability.} Detection rerun with two other random seeds and
tolerance factors 1.25 and 1.75; for each fault, the median Jaccard
overlap of its events with the best match; ``stable'' if \ensuremath{\geq} 0.5 (50 of
141). Overlap is low where parallel planes trade events between runs,
even when the geometry barely moves.

\textbf{Blanco strands} (\texttt{workflow/07\_faults/split\_strands.py},
\texttt{merge\_surfaces.py}). Strands are lines in map view within 20\textdegree{}
of plate motion (111\textdegree{}), 1.5 km half width, grown across gaps under 4 km,
merged when collinear. Two overlapping parallel strands stay separate
only if Ashman's D \textgreater{} 2 \citep{Ashman1994} over the overlap
and within each year with \ensuremath{\geq} 15 events in each. Depth test: the dip a
single fault would need to produce the offset. The literature check then
removed six NW-trending strands in the Cascadia Depression (where
published faults strike N--S) and merged two depth-dependent ``deeper
band'' strands. Table S2
(\texttt{fault\_surfaces\_v0\_strand\_separation.csv}) lists the 22
remaining pairs. Fig. S2: Blanco strands with their events and top
edges.

\section{Text S3. Removed surfaces}\label{text-s3.-removed-surfaces}

Six surfaces were quarry blasts (mostly daytime, weekday events: ON012,
ON020, ON023, ON029, ON037, ON040), four were artefacts (BL022, BL024,
BL026: deep planes where published Blanco seismicity stops at 10--14 km;
PU010: a seismogenic-layer artefact), six were strand-search artefacts
and two were merged (\texttt{fault\_surfaces\_v0\_removed.csv}, with
reasons).

\section{Text S4. Comparison with the
literature}\label{text-s4.-comparison-with-the-literature}

Five searches, one per region (Blanco; Mendocino triple junction; Gorda;
Puget and offshore Washington; onshore, offshore Oregon and Endeavour),
compared each fault's position, strike, dip, length and depth with
published faults, ridge and transform segments, seismic zones and
earthquake sequences. Matches carry a confidence (high to none) and a
class (known structure, duplicate, seismicity layer, no known
counterpart). Every DOI was checked on Crossref or DataCite; where no
structure is known, the references describe the area so the claim can be
checked. USGS moment tensors were used only to compare surfaces with
known earthquakes, never to build them. Table S4:
\texttt{fault\_surfaces\_v0\_literature.csv} and
\texttt{fault\_surfaces\_v0\_literature.html}; regional notes in
\texttt{data/faults/v0/literature/}.

\section{Text S5. Building the
catalog}\label{text-s5.-building-the-catalog}

(Adapted from the catalog draft, \texttt{paper/main.qmd}, Data and
Methods.)

\textbf{Data.} Continuous data from 2010 to 2016 were retrieved with
ROVER into a day-file archive read through an indexed client (pnwstore);
the \texttt{NC} and \texttt{BK} networks were requested from the NCEDC
FDSN service. Picks of catalog v4.0 come from 696 stations of 24
networks (\texttt{7A}, \texttt{7D}, \texttt{BK}, \texttt{C8},
\texttt{CC}, \texttt{CN}, \texttt{GS}, \texttt{NC}, \texttt{NV},
\texttt{OO}, \texttt{PB}, \texttt{TA}, \texttt{UO}, \texttt{UW},
\texttt{X4}, \texttt{X9}, \texttt{XD}, \texttt{XG}, \texttt{XU},
\texttt{YG}, \texttt{YH}, \texttt{YW}, \texttt{Z5}, \texttt{ZZ}); the
2016 repicks of v4.1 add 105 land stations (71 \texttt{NC}, 13
\texttt{NP}, 7 \texttt{BK}, 6 \texttt{UW}, 4 \texttt{CE}, 2 \texttt{YK},
1 \texttt{UO}, 1 \texttt{PB}), among them the Cascadia Initiative
ocean-bottom stations \citep{Stone2018} and the cabled stations of the
Ocean Observatories Initiative at Axial Seamount.

\textbf{Picking.} ELEP \citep{Yuan2023} combines five EQTransformer
models \citep{Mousavi2020} distributed by SeisBench \citep{Woollam2022}
(original, ETHZ, INSTANCE, SCEDC, STEAD) by second-order semblance over
0.5 s, on 60 s windows stepped by 30 s, after a 4--15 Hz band-pass; a
pick is the semblance maximum of each segment above 0.05. Two region
passes (2011--2015, 123--127\textdegree{}W; 2010--2015, 122--129\textdegree{}W) produced the
ver3 table (39.6 million picks). Three fills completed it: ver4 adds
region-run station-days left out of ver3 (9.9 million picks); ver5 adds
ocean-bottom station-days without picks, including the Axial cabled
stations (11.4 million); ver6 adds the land strip 122--123\textdegree{}W (12.7
million), for 73.6 million picks on 509,249 station-days. Every pick
keeps its identifier across versions.

\textbf{Association and location.} GENIE \citep{McBrearty2023}, trained
for 38--51\textdegree{}N, 119--132\textdegree{}W on synthetic events in a 3D model, associated
the ver6 picks day by day on the stations available each day; events
with at least five stations were kept and located, then relocated with
GraphDD \citep{McBrearty2025} on catalog and cross-correlation
differential times. Catalog v4 (I. McBrearty) has 172,340 events and
2,503,189 arrivals; v4.1, used here, has 172,342 events and 2,503,627
arrivals after the 2016 repick below
(\texttt{data/catalog\_v4/FIXES.md}). Quality metrics are the numbers of
P and S picks, stations, the azimuthal gap and the RMS residual; the
``QC'' subset used for fault surfaces requires at least 4 P and 4 S
picks.

\textbf{Magnitudes (Route A).} Wood--Anderson amplitudes are measured
once per station-day for every pick, in windows that grow with distance
after association (P: up to 1 + 0.03r s, S: up to 2 + 0.06r s); 99.6\%
of v4 arrivals have one. Station magnitudes with the Hutton--Boore
distance term are inverted for event and station terms (separate for
each deployment epoch) within 150 km, anchored to 3,416 ComCat local
magnitudes (scatter 0.21), and replaced by regional moment magnitudes
above M 4.5. 130,767 events have a magnitude.

\textbf{Known data issues.} Ocean-bottom records carry ship noise,
short-duration events, T-phases and high microseism noise on the shelf
\citep{Hilmo2020}; fragmented station-days produce picks at every data
join (1,846 station-days removed in ver6, 6,073 flagged in ver5);
\texttt{NC} and \texttt{BK} picks are missing in 2016 and parts of
2010--2011 (Text S7). Depth limitations: see below.

\textbf{Offshore location check.} On the Gorda Ridge, catalog epicenters
fall on the axial valley, while ComCat epicenters of the same events lie
10--20 km east.

Large events missed: 15 of 61 ComCat M \ensuremath{\geq} 5 (2010--2016) have no catalog
event within 30 s and 80 km, including the 8 December 2016 Mw 6.6 on the
Mendocino transform (10 at the northern edge, 3 offshore in 2016 without
ocean-bottom stations, 1 in the 2016 northern California pick gap, 1 on
the Blanco in 2011; script in the figure notes). Event 172113 (14:51:18,
4 P and 1 S, near Mount St.~Helens) is the Mw 6.6 mis-associated: its P
picks are the Pn arrivals at 708--797 km, within 2 s of iasp91; it is
removed in v4.1. We repicked the Mw 6.6, its M 4.7 aftershock and the
January 2016 Gorda pair on 30-minute windows of all stations within 4\textdegree{}
(279 for the Mw 6.6, 193 of them \texttt{NC}/\texttt{BK}) with the same
ELEP ensemble, kept the picks near the predicted arrivals from the
ComCat origins, and located them with NonLinLoc in the 3D model of the
relocation and in iasp91. The picks are clean to 400 km (Mw 6.6: 146 P,
54 S), but every station is east of the events (nearest 156 km, gap
255--310\textdegree{}): depths go to the bottom of every search range and epicenters
move 7--41 km west of ComCat. Catalog v4.1 adds the three missing events
(IDs 900001--900003) with the ComCat hypocenter and magnitude and our
picks, and adds 160 repicks to the M 4.7 aftershock (event 172114),
which keeps its v4 hypocenter
(\texttt{data/catalog\_v4/fixes/m66\_2016/},
\texttt{utils/fix\_catalog\_v4.py}). Depth limitations: 23.5\% of events
above the surface with depths read below sea level, 16.2\% below the
local surface; 3D relocation in progress.

\section{Text S6. Clustering: validation and
sensitivity}\label{text-s6.-clustering-validation-and-sensitivity}

(Code: \texttt{workflow/06\_analysis/v4\_clusters\_zbz.py}; report
\texttt{data/catalog\_v4/clusters\_zbz/catalog\_v4\_clusters\_report.html}.)
- Mc per region and year (maximum curvature + 0.2); b by maximum
likelihood per region; \(d_f\) = 1.6 (measured epicentral D2 0.8--1.4,
run as a variant); epicentral distance floored at 0.5 km. - Parent
search over the whole catalog with a rupture-length cap; events below Mc
or without magnitude attached afterwards to their nearest parent. -
Shuffled catalogs (times permuted in each region): bimodality index
drops from 1.6--3.4 to 0.3--1.6; only 3--17\% of shuffled events pass
the real threshold (34\% at Axial, whose clustered share is an upper
bound). - Validation: the 2010 and 2014 Gorda, 2011 and 2014 Vancouver
Island, 2012 Explorer and 2015 Mendocino fault zone earthquakes are
mainshock--aftershock sequences; the 2015 Axial eruption is one swarm. -
Sensitivity: 84--99\% of events keep their background or clustered label
under \(d_f\) 1.3/2.0/estimated, 3D distance, Mc + 0.3 and yearly Mc; a
conservative threshold raises background by 1--32 points. - Limits: the
per-region Mc keeps 27,555 events; Mc differs between regions
(0.75--2.45); the swarm rule's magnitude gap (0.5) is close to the ML
uncertainty (about 0.2).

\section{Text S7. Magnitude of completeness
maps}\label{text-s7.-magnitude-of-completeness-maps}

(Code: \texttt{workflow/06\_analysis/v4\_mc\_maps.py}; README in
\texttt{data/catalog\_v4/mc\_maps/}.) Magnitudes rounded to 0.1; every b
fit needs \ensuremath{\geq} 30 events above the cutoff. Methods: maximum curvature +0.2;
b-value stability with the Shi--Bolt test over a five-bin window
\citep{Woessner2005, Shi1982} (the fixed 0.03 tolerance of the
AgentSeismo code almost never passes with 250 events); goodness of fit
on cumulative counts at 90\% and 95\% \citep{Wiemer2000}. Maximum
curvature + 0.2 is the most conservative at 84\% of nodes; b-value
stability averages 0.19 lower and the best estimate 0.14 lower. The
QC-only catalog (\ensuremath{\geq} 4 P and \ensuremath{\geq} 4 S) leaves Mc unchanged (median difference
0.0, interquartile range \ensuremath{\pm}0.1). The per-region Mc of the clustering is
at or above all three methods everywhere except the Blanco
b-value-stability value (2.5 against 2.45). At Axial the
frequency--magnitude distribution has two slopes, so the pooled
goodness-of-fit value is unstable (\ensuremath{\sigma} 0.94); use maximum curvature (0.7)
or the node values (0.4--1.0).

\section{Text S8. Location versions and their
differences}\label{text-s8.-location-versions-and-their-differences}

(Code: \texttt{workflow/11\_compare/compare\_catalogs.py} and
\texttt{fig\_compare.py}; registry
\texttt{workflow/11\_compare/catalogs.yaml}; outputs and
\texttt{SUMMARY.md} in \texttt{data/catalog\_compare/}.) Each location
run is entered once in the registry with its velocity model, method,
parent catalog and status, and is read into one table layout. Events are
linked to the catalog v4 event IDs, directly for runs that relocate v4
events and by origin time (\ensuremath{\leq} 3 s, \ensuremath{\leq} 50 km) for catalog v3. The same
measures are computed for every version, overall and in eleven regions:
depth percentiles (km below sea level), the share of events above the
topography or bathymetry (GMT earth\_relief, 1 arc-minute), the shift
from catalog v4.1 for the same events, and the depth and epicenter
differences with ComCat events matched to v4.1 (16,584 events, \ensuremath{\leq} 3 s and
\ensuremath{\leq} 50 km). ``Near'' events have a station with a pick closer than
max(depth, 5 km). The spread of events about each fault surface (RMS
distance to the best-fitting plane of its members) is also computed, but
it favors v4, on which the surfaces were fitted, and it is lower for
double-difference catalogs than for single-event locations by
construction.

Current versions (Table S7, Fig. S8): v3 (Zenodo record; ver3 picks,
GraphDD), v4.0 and v4.1 (GENIE and GraphDD on ver6 picks in 1D models;
v4.1 adds the 2016 repick), and the NonLinLoc relocation of the v4
events in a 3D starting model without (iteration 0) and with (iteration
1) static station terms. Iteration 1 locates 157,267 events; the 15,073
others keep their v4 location (tier U). Station terms change locations
little (median 0.8 km horizontally, 0.1 km in depth). Against v4.1,
iteration 1 moves epicenters by a median of 2.4 km and depths by +2.7 km
(10--90\%: \ensuremath{-}3.7 to +10.6 km). Against ComCat, near events go from \ensuremath{-}1.0 \ensuremath{\pm}
2.4 km (v4.1) to +0.4 \ensuremath{\pm} 2.4 km (median \ensuremath{\pm} MAD); all matched events from
\ensuremath{-}0.9 \ensuremath{\pm} 3.8 km to +1.2 \ensuremath{\pm} 4.2 km, the deep bias of poorly constrained
events. By region, the 3D locations remove the shallow bias of v4 at
Axial (median 0.3 to 3.0 km, 71\% to 0\% above the seafloor) and on the
Juan de Fuca ridges, and push Gorda (+3.7 km against ComCat, near
events) and Mendocino (+2.8 km) deeper. Each new run (tomography
relocations, GrowClust) is added to the registry and to this text.

\section{Supplementary tables (files)}\label{supplementary-tables-files}

\begin{itemize}
\tightlist
\item
  Table S1. Events and share on a fault, per region.
\item
  Table S2. Blanco strand separation tests.
\item
  Table S3. Fault surfaces: geometry, uncertainty, stability
  (\texttt{fault\_surfaces\_v0\_summary.csv}).
\item
  Table S4. Literature matches with verified references.
\item
  Table S5. Clustering families (\texttt{families\_zbz.csv}) and
  per-region summary (\texttt{regions\_zbz.csv}).
\item
  Table S6. Magnitude of completeness by region and period, three
  methods, with bootstrap sigma, b and rate ratios
  (\texttt{data/catalog\_v4/mc\_maps/mc\_regional.csv}); grids in
  \texttt{mc\_grid.csv},
  \texttt{mc\_grid\_\textless{}period\textgreater{}.csv} and NetCDF.
  Figs. S3--S7: three-method maps, bootstrap sigma and b, periods,
  regional frequency--magnitude distributions, QC variant
  (\texttt{data/catalog\_v4/mc\_maps/fig1--fig5}).
\item
  Table S7. Location versions compared: whole catalog and per region
  (\texttt{data/catalog\_compare/summary.csv}, \texttt{anss.csv},
  \texttt{pairs.csv}, \texttt{faults\_summary.csv}; readable tables in
  \texttt{SUMMARY.md}). Fig. S8: depth, depth against ComCat and events
  above the surface by region and version, and the change from v4.1 to
  the latest version (\texttt{figures/figS\_catalog\_compare}).
\end{itemize}

\bibliographystyle{unsrtnat}
\bibliography{references}
\end{document}
